\section{Runtime choices and compiled implementations}
\label{sec:arch-algorithm-selection}

An alignment request names biological work, but it does not identify one
machine-code function. BioSeqDB must also account for the sequence alphabet,
packed representation, available indexes, and execution hardware. These
decisions select an already compiled implementation and connect the algorithms
described earlier to the storage and execution layers.

The logical flow remains sketching, seed lookup, chaining, region selection,
and optional base alignment. Sections~\ref{sec:minimizers-role},
\ref{sec:chaining-bioseqdb}, and~\ref{sec:final-alignment-role} explain those
algorithms. Here the concern is how the system selects their implementations
and supplies compatible inputs.

\subsection{Turn SQL settings into a request}
\label{sec:arch-request-settings}

PostgreSQL configuration parameters, usually called GUCs, provide session
defaults. The three settings in Listing~\ref{lst:arch-algorithm-options}
select minimizer extraction, chaining and base alignment, respectively.
Their defaults are \texttt{scalar}, \texttt{scalar\_minimap2} and
\texttt{block\_aligner}. The available CPU minimizer choices are described in
Section~\ref{sec:configured-cpu-minimizers} and
Table~\ref{tab:current-cpu-minimizers}.

A constant \texttt{aligner\_options} record overrides individual settings
for one search. Unspecified or NULL fields inherit the session values.
The following example instead selects all three algorithm families through
named GUC values: scalar sketching, scalar chaining, and KSW2 base alignment.
Both \texttt{refs} and \texttt{queries} are stored sequence tables with an
\texttt{id} column and a \texttt{seq} column of type
\texttt{bioseq.seq('dna4',32)}.

\noindent\begin{minipage}{\linewidth}
\begin{lstlisting}[style=bioseq-search-sql,
  caption={Select algorithm families before planning a sequence join.},
  label={lst:arch-algorithm-options}]
BEGIN;
SET LOCAL bioseq.aligner.minimizer.algorithm = 'scalar';
SET LOCAL bioseq.aligner.chaining.algorithm = 'scalar_minimap2';
SET LOCAL bioseq.aligner.base_alignment.algorithm = 'ksw2';
EXPLAIN
SELECT r.id, q.id
FROM refs AS r JOIN queries AS q ON mmsearch(r.seq, q.seq);
COMMIT;
\end{lstlisting}
\end{minipage}

An ordinary \texttt{BioseqScan} stores the defaults and overrides in its
plan. A \texttt{BioseqBatchJoin} reads the defaults again when execution
starts and reapplies the constant overrides. Changes to the corresponding
alignment GUCs invalidate cached plans.
The request also fixes the selected chaining policy:
PostgreSQL SIMD chaining uses \texttt{bounded\_\allowbreak batch\_\allowbreak approximate}, whereas
the C++ interface additionally permits \texttt{exact\_visited}. They are
distinct policies, as described in Section~\ref{sec:chaining-bioseqdb}.
Likewise, the internal batched KSW2 implementation is not an additional
PostgreSQL base-aligner setting.

Required results also determine which stages run. A coordinate-only request
can avoid base alignment. A CIGAR or exact edit-distance request requires
additional alignment work. Divergence and mapping quality have their own
requirements. Thus selecting an aligner does not imply that every search
executes it. A prebuilt query sketch also lets execution enter after
sketching.

\subsection{Separate algorithm selection from hardware selection}
\label{sec:arch-hardware-selection}

Figure~\ref{fig:arch-selection} separates the decisions that constrain
execution. A requested combination must have an implementation and satisfy
its input contract. The system does not treat every combination of algorithm
names as interchangeable or supported.

\begin{figure}[htbp]
\centering
\begin{tikzpicture}[
  x=1cm,y=1cm,font=\footnotesize,
  arch input/.style={draw=black!55,fill=black!3,rounded corners=2pt,
    minimum height=.88cm,text width=3.55cm,align=center,inner sep=3pt},
  arch decision/.style={draw=blue!60!black,fill=blue!7,rounded corners=2pt,
    minimum height=1.00cm,text width=4.35cm,align=center,inner sep=4pt},
  arch flow/.style={-{Latex[length=1.8mm]},thick}
]
  \node[arch input,text width=4.35cm] (request) at (6.8,4.25)
    {Requested algorithms\\and required results};
  \node[arch input] (compiled) at (2.05,2.65)
    {Precompiled stage variants\\and representation choices};
  \node[arch decision] (admit) at (6.8,2.65)
    {Select a supported path\\and check its input contract};
  \node[arch input] (inputs) at (11.55,2.65)
    {Alphabet, logical ranges,\\and stored-index identity};
  \node[arch input] (cpu) at (2.05,.85)
    {CPU paths: instruction set\\compiled and CPU-supported};
  \node[arch decision,fill=green!7,draw=green!45!black] (execute) at (6.8,.85)
    {Execute the admitted path\\with its selected kernels};
  \node[arch input,fill=red!5,draw=red!50!black] (reject) at (11.55,.85)
    {Reject an unsupported\\request};
  \draw[arch flow] (request)--(admit);
  \draw[arch flow] (compiled)--(admit);
  \draw[arch flow] (inputs)--(admit);
  \draw[arch flow] (cpu)--(execute);
  \draw[arch flow] (admit)--(execute);
  \draw[arch flow] (admit.south east)--(reject.north west);
\end{tikzpicture}
\caption{Algorithm selection, input compatibility, and CPU target selection
have different responsibilities. The diagram shows their dependencies,
not a universal order of planner callbacks. GPU execution has additional
admission requirements. A runtime failure can still occur after admission.}
\label{fig:arch-selection}
\end{figure}

For CPU kernels, Highway selects an instruction-set target that is both
compiled into the extension and supported by the running processor.
BioSeqDB binds the resulting kernel table once at the alignment-batch entry
and passes it to the selected implementation. It does not repeat CPU feature
detection for each seed or choose a different biological algorithm because
another instruction set is available. PostgreSQL vectorized tuple processing
is a separate execution choice from the SIMD kernels inside an aligner.

Here a \emph{specialization} is a compiled variant of a C++ template.
Runtime selectors choose among these variants. The dispatch chain fixes
a representation bucket and algorithm traits before
calling a typed pipeline. Representation-width choices depend on the
selected k-mer bucket, avoiding a full independent product of both axes.
When an outer caller has already selected the representation, the inner
pipeline retains that choice. When sketches are supplied, it also removes
the unnecessary sketch-algorithm dispatch dimension.

Not every numerical parameter needs its own specialization. Compile profiles
control which specialized forms are instantiated. Other values use a compiled
generic implementation that reads the requested $k$ and $w$ at runtime. This
\emph{specialization fallback} preserves the biological parameters. It is
different from the explicit short-pair adjustment of untouched defaults,
which can choose smaller $k$ and $w$ before execution. An indexed path
prepares query sketches with the parameters required by the stored lookup.

\subsection{Require compatible indexes and execution routes}
\label{sec:arch-route-admission}

An index is usable only when its identity matches the selected operation.
The relevant information includes the source generation, alphabet,
minimizer scheme, producer identity, and $k,w$ values. Matching the latter
two numbers alone is insufficient. In particular, the resident and indexed
continuation routes within the two-stack family cannot be mixed merely
because they share an algorithm name. The indexed query uses the selection
route required by its index. Section~\ref{sec:index-runtime-selection}
describes the separate policies for allowing or requiring persisted reuse.

Admission also limits acceleration:
\begin{itemize}
  \item \textbf{CPU SIMD.} Selected kernels check their supported alphabet,
    representation, parameter range, and vector shape. Bounded ntHash
    has registered DNA4 and DNA5 kernels with both specialized and generic
    parameter handling, as summarized in
    Section~\ref{sec:configured-cpu-minimizers}. SIMD chaining requires a vector target
    and valid limits for its selected policy. A rejected request is not
    retried with another chaining policy or a scalar implementation.
  \item \textbf{Biological Vulkan stages.} PostgreSQL requires the complete
    supported resident combination, including Vulkan sketching, seed lookup,
    chaining, fused execution, and permission to build GPU indexes when
    needed. It requires
    usable hardware and a batch containing at least two jobs. Unsupported
    mixtures and device failures produce errors instead of entering the
    general CPU pipeline.
  \item \textbf{GPU output.} The persisted grouped-SoA route currently
    supports coordinates only. Transient resident execution also has a
    bounded CIGAR route, which retains host packing and result construction.
    Admission to a GPU route therefore does not mean that every substage
    runs on the device.
\end{itemize}

These rules have specific exceptions and should not be reduced to a blanket
ban on fallback. Historical reference-SIMD measurements can explicitly record
scalar execution, whereas current registered DNA5 profiles have a SIMD route.
A Vulkan rank prefilter
can accompany CPU biological stages. A permissive index policy can rebuild
an absent or inapplicable lookup. Finally, a rejected parallel work split can
run serially with the same algorithm. Each alternative belongs to a
particular layer and preserves that layer's stated contract.
